\documentclass[12pt,oneside]{book}

% ============================================================
%  PACKAGES
% ============================================================
\usepackage[T1]{fontenc}
\usepackage[utf8]{inputenc}
\usepackage[english]{babel}
\usepackage{geometry}
\geometry{
  paperwidth=7in,
  paperheight=9.25in,
  top=1in,
  bottom=1in,
  inner=1.25in,
  outer=0.75in
}

% Typography
\usepackage{lmodern}
\usepackage{microtype}
\usepackage{setspace}
\onehalfspacing

% Code listings
\usepackage{listings}
\usepackage{xcolor}

\definecolor{codebg}{RGB}{245,245,245}
\definecolor{codeframe}{RGB}{200,200,200}
\definecolor{keyword}{RGB}{0,0,180}
\definecolor{string}{RGB}{163,21,21}
\definecolor{comment}{RGB}{80,130,80}

\lstdefinestyle{pythonstyle}{
  backgroundcolor=\color{codebg},
  commentstyle=\color{comment}\itshape,
  keywordstyle=\color{keyword}\bfseries,
  stringstyle=\color{string},
  basicstyle=\ttfamily\footnotesize,
  breakatwhitespace=false,
  breaklines=true,
  captionpos=b,
  keepspaces=true,
  numberstyle=\tiny\color{gray},
  numbers=left,
  numbersep=8pt,
  showspaces=false,
  showstringspaces=false,
  showtabs=false,
  tabsize=4,
  frame=single,
  rulecolor=\color{codeframe},
  language=Python
}
\lstset{style=pythonstyle}

% Headers & navigation
\usepackage{fancyhdr}
\pagestyle{fancy}
\fancyhf{}
\fancyhead[L]{\leftmark}
\fancyhead[R]{\thepage}
\renewcommand{\headrulewidth}{0.4pt}

% Hyperlinks
\usepackage[colorlinks=true,linkcolor=blue,urlcolor=blue,citecolor=blue]{hyperref}

% Callout boxes (tips, notes, warnings)
\usepackage{tcolorbox}
\tcbuselibrary{skins,breakable}
\newtcolorbox{notebox}[1][Note]{
  colback=blue!5!white, colframe=blue!40!black,
  title={\textbf{#1}}, breakable
}
\newtcolorbox{warnbox}[1][Warning]{
  colback=red!5!white, colframe=red!50!black,
  title={\textbf{#1}}, breakable
}
\newtcolorbox{tipbox}[1][Tip]{
  colback=green!5!white, colframe=green!50!black,
  title={\textbf{#1}}, breakable
}

% Graphics
\usepackage{graphicx}
\usepackage{float}

% Tables
\usepackage{booktabs}
\usepackage{tabularx}

% Index
\usepackage{makeidx}
\makeindex

% ============================================================
%  METADATA
% ============================================================
\title{
  {\Huge \textbf{Building FrenchBooks}}\\[0.5em]
  {\large A Complete Guide to Building a Native Document Reader\\
  with PyQt6, spaCy, and Python}\\[1em]
  {\normalsize DO-178C Traceability Edition}
}
\author{AsterionDantes}
\date{2026-06-30 \\ Version 1.0.0}

% ============================================================
\begin{document}
% ============================================================

\maketitle
\thispagestyle{empty}
\newpage

% ============================================================
\frontmatter
% ============================================================

\tableofcontents
\newpage

% ============================================================
%  PREFACE
% ============================================================
\chapter*{Preface}
\addcontentsline{toc}{chapter}{Preface}

This book explains, from first principles, every technical decision made in building
\textbf{FrenchBooks}—an interactive French literature study tool that opens PDF and EPUB
books natively, processes them with a French NLP engine, and dynamically generates a
per-page study guide.

No prior programming experience is assumed. If you have picked up this document, you are
either curious how a modern desktop application is assembled from its components, or you
need to recreate, maintain, or certify this software. Either way, this book will walk you
through every concept, every library, and every line of code.

The project is built to \textbf{DO-178C} standards—a software assurance framework
originally developed for aircraft avionics software. Applying these standards to a desktop
application is deliberate: it enforces a discipline of traceability, documentation, and
auditability that benefits any serious software project, not just aerospace systems.

\textbf{How this book is organized:}

\begin{itemize}
  \item \textbf{Part I — Foundations} covers the tools and concepts you need before
    writing a single line of code: Python virtual environments, the PyQt6 event-driven
    programming model, and what ``MVC'' actually means.
  \item \textbf{Part II — The Reader} covers building the split-screen document viewer:
    rendering PDF pages as images and rendering EPUB HTML in a web engine widget.
  \item \textbf{Part III — The NLP Engine} covers the French language processing pipeline:
    how spaCy understands French verbs, tenses, and lemmas, and how to run all of that
    without freezing the user interface.
  \item \textbf{Part IV — The Dashboard} covers corpus-level analysis: word frequency
    heatmaps, tense summaries, and vocabulary lists.
  \item \textbf{Part V — Certification} covers the DO-178C traceability pipeline, GitFlow,
    and how to maintain the audit trail.
\end{itemize}

% ============================================================
\mainmatter
% ============================================================

% ============================================================
%  PART I — FOUNDATIONS
% ============================================================
\part{Foundations}

% ============================================================
\chapter{Introduction to the Project}
\label{ch:intro}
% ============================================================

\section{What is FrenchBooks?}

FrenchBooks is a \textbf{desktop application}—a program that runs directly on your
computer, rather than inside a web browser. When you open it, you see a window with
panels, buttons, and scrollable text. This is different from a website (which lives on a
server) or a mobile app (which lives on a phone).

The application has one job: to make studying French literature as efficient as possible.
It does this by reading a French book (either a PDF or an EPUB file), and while you read
it, it simultaneously runs a linguistic analysis of every word on the page and presents
you with a live study guide on the screen next to the book.

\section{The Technology Stack}

A ``technology stack'' is the set of software tools that work together to make an
application function. Think of it as the list of ingredients in a recipe. Table
\ref{tab:stack} lists ours.

\begin{table}[H]
\centering
\caption{FrenchBooks Technology Stack}
\label{tab:stack}
\begin{tabularx}{\textwidth}{l X}
\toprule
\textbf{Tool} & \textbf{Role in the Application} \\
\midrule
Python 3.11 & The programming language. All application logic is written in Python. \\
PyQt6 & The GUI toolkit. Provides windows, buttons, layouts, and the event loop. \\
PyQt6-WebEngine & A web browser widget embedded inside PyQt6. Used to render EPUB HTML. \\
PyMuPDF (fitz) & A PDF library. Converts PDF pages into images for display. \\
EbookLib & An EPUB library. Opens and reads the structure of EPUB books. \\
spaCy & A Natural Language Processing library. Understands French grammar. \\
fr\_core\_news\_md & The French language model for spaCy. Contains learned French grammar rules. \\
matplotlib & A charting library. Draws the word frequency heatmap. \\
BeautifulSoup4 & An HTML parsing library. Strips HTML tags from EPUB content for NLP. \\
\bottomrule
\end{tabularx}
\end{table}

\section{The DO-178C Standard}

DO-178C (full title: \textit{Software Considerations in Airborne Systems and Equipment
Certification}) is a set of guidelines published by RTCA, Inc. Originally written for
software in aircraft control systems—where a software bug can be catastrophic—it defines
a rigorous process for:

\begin{enumerate}
  \item Writing requirements before writing code
  \item Ensuring every line of code can be traced back to a requirement
  \item Maintaining comprehensive documentation that would allow a stranger to
    understand, verify, and recreate the software
\end{enumerate}

We apply these principles here not because FrenchBooks will fly an airplane, but because
they enforce good engineering discipline. The \texttt{Requirements.md} file is our
``Software Requirements Specification.'' This book (\texttt{Documentation.tex}) is our
``Software Design Document.'' The \texttt{Recreate.md} file is our
``Configuration Management Log.''

% ============================================================
\chapter{Python Virtual Environments}
\label{ch:venv}
% ============================================================

\section{The Problem: Conflicting Dependencies}

Imagine you are working on two Python projects simultaneously. Project A needs version 1.0
of a library called ``colorize.'' Project B needs version 2.0 of the same library—and
version 2.0 has a completely different interface that is incompatible with version 1.0.

If you install both on the same Python installation, one project will always break. This
is the ``dependency conflict'' problem.

\section{The Solution: Virtual Environments}

A \textbf{virtual environment} (often abbreviated ``venv'') is a self-contained copy of
Python, isolated from the rest of your system. When you activate a virtual environment:

\begin{itemize}
  \item The \texttt{python} and \texttt{pip} commands point to a private copy inside the
    venv folder.
  \item Any packages you install go only into the venv, not into the system Python.
  \item Deleting the venv folder removes all its packages cleanly, without affecting
    anything else on your machine.
\end{itemize}

\section{Creating and Activating a venv (Windows PowerShell)}

\begin{lstlisting}[language=bash, caption={Virtual environment commands}]
# Step 1: Create the virtual environment
python -m venv venv

# Step 2: Activate it
.\venv\Scripts\Activate.ps1

# Your prompt now shows (venv) at the left
# Step 3: Install packages (they go into venv only)
pip install PyQt6

# Step 4: Deactivate when you are done
deactivate
\end{lstlisting}

\begin{notebox}
The \texttt{venv/} folder is excluded from git via \texttt{.gitignore}. It is never
committed to the repository. Anyone who clones the repo recreates it from scratch using
\texttt{Recreate.md}.
\end{notebox}

% ============================================================
\chapter{The Model-View-Controller Pattern}
\label{ch:mvc}
% ============================================================

\section{The Problem: Spaghetti Code}

In a small application, it is tempting to put everything in one file. The button click
handler reads the file, parses it, runs NLP on it, formats the results, and updates the
display—all in one function. This works until the application grows. Soon, changing one
thing breaks another, and nobody can understand where anything is.

\section{The MVC Solution}

The \textbf{Model-View-Controller} (MVC) pattern divides an application into three
separate roles:

\begin{description}
  \item[Model] \hfill \\
    The \textit{data brain}. Knows what data exists, how to load it, and how to transform
    it. Has no idea how the screen looks. In FrenchBooks, the Model opens PDF/EPUB files,
    caches text, and stores NLP results.

  \item[View] \hfill \\
    The \textit{face}. Knows how to display data beautifully. Has no idea where the data
    came from. In FrenchBooks, the View contains PyQt6 widgets—scroll areas, tables,
    labels—that show whatever data is handed to them.

  \item[Controller] \hfill \\
    The \textit{traffic director}. Listens for user actions (button clicks, page turns),
    asks the Model for data, and tells the View what to display. In FrenchBooks, the
    Controller dispatches NLP worker threads when the user turns a page.
\end{description}

\begin{figure}[H]
\centering
% Diagram placeholder — replace with actual TikZ diagram in later revision
\fbox{\parbox{0.6\textwidth}{\centering
  \textbf{[MVC Interaction Diagram]} \\[1em]
  User $\rightarrow$ View (click event) \\
  View $\rightarrow$ Controller (signal) \\
  Controller $\rightarrow$ Model (fetch data) \\
  Controller $\rightarrow$ NLPWorker (dispatch thread) \\
  NLPWorker $\rightarrow$ Controller (finished signal) \\
  Controller $\rightarrow$ View (update display)
}}
\caption{MVC interaction flow in FrenchBooks}
\label{fig:mvc}
\end{figure}

\section{Why MVC Matters for DO-178C}

DO-178C requires that every piece of software can be tested independently. If your
NLP logic is tangled inside your display code, you cannot test them separately. With MVC,
you can write unit tests for the Model (does it correctly parse a PDF?) entirely
independently of the View (how does the table look?).

% ============================================================
\chapter{PyQt6 and Event-Driven Programming}
\label{ch:pyqt6}
% ============================================================

\section{What is a GUI Framework?}

A \textbf{GUI framework} is a library that provides the building blocks for a graphical
user interface: windows, buttons, text boxes, menus, scroll bars, and so on. More
importantly, it provides the \textit{event loop}—the engine that keeps the application
alive, listening for user input, and responding to it.

PyQt6 is the Python binding for Qt 6, a professional C++ GUI toolkit used in commercial
applications worldwide (including many embedded systems and, yes, aerospace HMIs).

\section{Signals and Slots: Qt's Communication System}

Qt uses a communication mechanism called \textbf{signals and slots}. This is fundamentally
different from ordinary function calls.

\begin{description}
  \item[Signal] \hfill \\
    An event that a widget \textit{announces}. For example, a button has a
    \texttt{clicked} signal that fires whenever the user clicks it.
  \item[Slot] \hfill \\
    A function that \textit{listens} for a signal. You \textit{connect} a signal to a
    slot; when the signal fires, the slot function is called automatically.
\end{description}

\begin{lstlisting}[caption={Signal and slot example in PyQt6}]
from PyQt6.QtWidgets import QPushButton

button = QPushButton("Analyze Page")

# Connect the clicked signal to a slot (any callable)
button.clicked.connect(self.on_analyze_clicked)

def on_analyze_clicked(self):
    # This runs automatically when the button is clicked
    self.controller.dispatch_nlp_worker(current_page_text)
\end{lstlisting}

\section{The Application and Event Loop}

Every PyQt6 application must create exactly one \texttt{QApplication} object and call
\texttt{app.exec()} to start the event loop. The event loop runs continuously, processing
keyboard presses, mouse clicks, and timer events, until the user closes the last window.

\begin{lstlisting}[caption={Minimal PyQt6 application entry point}]
import sys
from PyQt6.QtWidgets import QApplication

def main():
    app = QApplication(sys.argv)
    # ... create and show windows ...
    sys.exit(app.exec())

if __name__ == "__main__":
    main()
\end{lstlisting}

% ============================================================
%  PART II — THE READER
% ============================================================
\part{The Native Document Reader}

% ============================================================
\chapter{Rendering PDF Pages with PyMuPDF}
\label{ch:pdf}
% ============================================================

\section{What is PyMuPDF?}

PyMuPDF (imported in Python as \texttt{fitz}) is a Python interface to MuPDF, a
lightweight, high-performance PDF and XPS rendering engine written in C. The name
``MuPDF'' comes from ``minimal'' PDF. It can open PDF files, parse their internal
structure, and convert pages to pixel data with high fidelity and speed.

\section{The Conceptual Pipeline: PDF Page to Qt Screen}

Displaying a PDF page in a PyQt6 window involves four distinct steps:

\begin{enumerate}
  \item \textbf{Open} the PDF file using PyMuPDF.
  \item \textbf{Render} a specific page into a \texttt{Pixmap}—a grid of raw pixel values
    (RGB data) representing the page as a bitmap image.
  \item \textbf{Convert} the raw pixel data into a \texttt{QImage} (a Qt object that
    understands image data).
  \item \textbf{Display} the \texttt{QImage} as a \texttt{QPixmap} inside a
    \texttt{QLabel} widget embedded in a scrollable area.
\end{enumerate}

\section{Step 1: Opening a PDF File}

\begin{lstlisting}[caption={Opening a PDF with PyMuPDF (LLR-04)}]
import fitz  # PyMuPDF

# Open the PDF file. Returns a Document object.
doc = fitz.open("path/to/book.pdf")

# How many pages does it have?
print(f"Page count: {doc.page_count}")
\end{lstlisting}

The \texttt{fitz.open()} function reads the PDF file into memory and returns a
\texttt{Document} object. This object is your handle to everything inside the PDF.

\section{Step 2: Rendering a Page to a Pixmap}

\begin{lstlisting}[caption={Rendering a page as a bitmap (LLR-04)}]
# Load page number 0 (PDF pages are zero-indexed)
page = doc.load_page(0)

# Create a transformation matrix for scaling.
# zoom=1.5 means 150% of 72 DPI = ~108 DPI.
zoom = 1.5
matrix = fitz.Matrix(zoom, zoom)

# Render the page. Returns a Pixmap (raw pixel grid).
pixmap = page.get_pixmap(matrix=matrix)
\end{lstlisting}

\begin{notebox}
\textbf{What is a Matrix?} In mathematics, a matrix is a grid of numbers used to
transform 2D coordinates. PyMuPDF's \texttt{Matrix(zoom, zoom)} scales the page up by
the zoom factor in both the X and Y directions. A zoom of 1.0 renders at 72 DPI (PDF
native resolution). A zoom of 2.0 renders at 144 DPI (sharper on high-DPI screens).
\end{notebox}

\section{Step 3: Converting Pixmap to QImage}

A PyMuPDF \texttt{Pixmap} is not directly understood by PyQt6. We must convert it:

\begin{lstlisting}[caption={Converting fitz.Pixmap to QImage (LLR-05)}]
from PyQt6.QtGui import QImage, QPixmap

# pixmap.samples is a bytes object: raw RGB pixel data
# pixmap.width, pixmap.height: pixel dimensions
# pixmap.stride: bytes per row (width * 3 for RGB)
image = QImage(
    pixmap.samples,
    pixmap.width,
    pixmap.height,
    pixmap.stride,
    QImage.Format.Format_RGB888   # 3 bytes per pixel: Red, Green, Blue
)

# Convert QImage to QPixmap (the type Qt widgets display)
qt_pixmap = QPixmap.fromImage(image)
\end{lstlisting}

\begin{notebox}
\textbf{Why two different classes?} \texttt{QImage} is for image \textit{manipulation}—it
gives you pixel-level access. \texttt{QPixmap} is for image \textit{display}—it is
optimized by the GPU driver for fast rendering on screen. Qt keeps them separate because
not all displays can render arbitrary image formats, but all displays can render QPixmap.
\end{notebox}

\section{Step 4: Displaying in a Scrollable Widget}

\begin{lstlisting}[caption={Displaying QPixmap in a scroll area (LLR-06)}]
from PyQt6.QtWidgets import QLabel, QScrollArea
from PyQt6.QtCore import Qt

# A QLabel can display an image (QPixmap)
label = QLabel()
label.setPixmap(qt_pixmap)
label.setAlignment(Qt.AlignmentFlag.AlignTop | Qt.AlignmentFlag.AlignLeft)

# A QScrollArea allows panning when the image is larger than the visible area
scroll_area = QScrollArea()
scroll_area.setWidget(label)
scroll_area.setWidgetResizable(False)
\end{lstlisting}

% ============================================================
\chapter{Rendering EPUB Content with EbookLib and QWebEngineView}
\label{ch:epub}
% ============================================================

\section{What is an EPUB File?}

An EPUB (Electronic Publication) file is not a single document—it is a ZIP archive
containing a collection of HTML files, CSS stylesheets, images, and metadata. Think of it
as a tiny website, zipped up. The ``spine'' of an EPUB defines the reading order of its
HTML files (chapters).

\section{Parsing an EPUB with EbookLib}

\begin{lstlisting}[caption={Opening an EPUB and iterating chapters (LLR-07)}]
import ebooklib
from ebooklib import epub

# Open the EPUB file
book = epub.read_epub("path/to/book.epub")

# Get all HTML documents in reading order (the spine)
chapters = list(book.get_items_of_type(ebooklib.ITEM_DOCUMENT))

print(f"Number of chapters: {len(chapters)}")

# Get the raw HTML content of the first chapter
html_content = chapters[0].get_content().decode("utf-8")
\end{lstlisting}

\section{Rendering HTML in QWebEngineView}

\texttt{QWebEngineView} embeds a full Chromium-based web browser engine inside a PyQt6
widget. It can render the same HTML and CSS that Google Chrome understands.

\begin{lstlisting}[caption={Rendering EPUB HTML in QWebEngineView (LLR-08)}]
from PyQt6.QtWebEngineWidgets import QWebEngineView

# Create the web view widget
web_view = QWebEngineView()

# Load HTML string directly (no network request needed)
web_view.setHtml(html_content)
\end{lstlisting}

\begin{warnbox}
\texttt{QWebEngineView} requires the \texttt{PyQt6-WebEngine} package to be installed
separately from PyQt6 itself. If the import fails, fall back to \texttt{QTextBrowser},
which renders a limited subset of HTML. See Section \ref{sec:textbrowser-fallback}.
\end{warnbox}

\section{Extracting Plain Text from EPUB HTML}
\label{sec:epub-text}

The NLP pipeline needs plain text, not HTML. We use BeautifulSoup4 to strip the HTML tags:

\begin{lstlisting}[caption={Stripping HTML tags for NLP input}]
from bs4 import BeautifulSoup

soup = BeautifulSoup(html_content, "lxml")
plain_text = soup.get_text(separator=" ", strip=True)
# plain_text is now clean French prose, ready for spaCy
\end{lstlisting}

\section{QTextBrowser Fallback}
\label{sec:textbrowser-fallback}

If \texttt{PyQt6-WebEngine} is not installed (LLR-09):

\begin{lstlisting}[caption={QTextBrowser fallback for EPUB rendering}]
from PyQt6.QtWidgets import QTextBrowser

text_browser = QTextBrowser()
text_browser.setHtml(html_content)
# QTextBrowser renders basic HTML: headings, paragraphs, bold, italic
# It does NOT support JavaScript or complex CSS
\end{lstlisting}

% ============================================================
%  PART III — THE NLP ENGINE
% ============================================================
\part{The French NLP Engine}

% ============================================================
\chapter{Threading: Keeping the GUI Responsive}
\label{ch:threading}
% ============================================================

\section{The Problem: The Frozen Window}

When a Python program runs a slow operation—like loading a 400-page PDF or running
NLP on 5,000 words—it is busy. If that operation runs on the \textit{main thread}
(the thread that also runs the GUI), the window becomes completely unresponsive. You
cannot click anything. You cannot move the window. On Windows, it shows as ``Not
Responding'' in the title bar.

This is unacceptable (and a violation of HLR-04).

\section{Threads: Doing Two Things at Once}

A \textbf{thread} is an independent sequence of instructions that runs concurrently with
other threads in the same program. Modern computers have multiple CPU cores, so threads
can literally run at the same time on different cores.

The \textbf{main thread} in a PyQt6 application has one sacred responsibility: it runs
the event loop that keeps the GUI alive. \textit{No slow operation should ever run on the
main thread.}

\section{Qt's Threading Tools: QRunnable and QThreadPool}

Qt provides a convenient system for background work:

\begin{description}
  \item[\texttt{QRunnable}] A lightweight task that can be run in a background thread.
    You subclass it and put your slow code inside the \texttt{run()} method.
  \item[\texttt{QThreadPool}] A pool of pre-created background threads, managed by Qt.
    You submit a \texttt{QRunnable} to the pool, and Qt assigns it to an idle thread.
  \item[\texttt{QObject} signals] The \textit{only} safe way to send a result from a
    background thread back to the main (GUI) thread. Qt's signal mechanism is
    thread-safe by design.
\end{description}

\section{The NLPWorker Pattern}

\begin{lstlisting}[caption={The NLPWorker class (LLR-10, LLR-11, LLR-12)}]
from PyQt6.QtCore import QRunnable, QObject, pyqtSignal

class WorkerSignals(QObject):
    """Signals must live on a QObject, not on a QRunnable."""
    finished = pyqtSignal(object)  # carries the spaCy Doc result
    error    = pyqtSignal(str)

class NLPWorker(QRunnable):
    """Runs spaCy NLP in a background thread."""

    def __init__(self, nlp_pipeline, text):
        super().__init__()
        self.nlp = nlp_pipeline
        self.text = text
        self.signals = WorkerSignals()
        self._cancelled = False

    def cancel(self):
        """Thread-safe cancellation flag (LLR-13)."""
        self._cancelled = True

    def run(self):
        """This method runs in a background thread — never on the GUI thread."""
        if self._cancelled:
            return
        try:
            doc = self.nlp(self.text)  # The slow spaCy call
            if not self._cancelled:
                self.signals.finished.emit(doc)
        except Exception as e:
            self.signals.error.emit(str(e))
\end{lstlisting}

\begin{lstlisting}[caption={Dispatching the worker from the Controller (LLR-12, LLR-13)}]
from PyQt6.QtCore import QThreadPool

class ReaderController:
    def __init__(self, nlp_pipeline):
        self.nlp = nlp_pipeline
        self._current_worker = None

    def on_page_changed(self, page_text):
        # Cancel any previous worker still running
        if self._current_worker:
            self._current_worker.cancel()

        # Create and dispatch a new worker
        worker = NLPWorker(self.nlp, page_text)
        worker.signals.finished.connect(self.on_nlp_complete)
        self._current_worker = worker
        QThreadPool.globalInstance().start(worker)

    def on_nlp_complete(self, spacy_doc):
        # This slot runs on the MAIN thread (Qt delivers cross-thread signals safely)
        self.study_guide_view.update(spacy_doc)
\end{lstlisting}

\begin{warnbox}
\textbf{Never pass PyQt6 widget references into a QRunnable.} Workers run in background
threads; accessing widgets from a background thread causes race conditions and crashes.
Results must travel back via Qt signals only.
\end{warnbox}

% ============================================================
\chapter{French NLP with spaCy}
\label{ch:nlp}
% ============================================================

\section{What is Natural Language Processing?}

\textbf{Natural Language Processing} (NLP) is the field of computer science concerned
with enabling computers to understand, analyze, and generate human language. A French
sentence like \textit{``Il avait mangé la pomme''} (He had eaten the apple) is,
to a computer, just a string of characters. NLP gives the computer the tools to
understand that \textit{avait mangé} is a verb in the \textit{Plus-que-parfait} tense
and that \textit{pomme} is a noun meaning ``apple.''

\section{Loading the French spaCy Model}

\begin{lstlisting}[caption={Loading the French model (LLR-12)}]
import spacy

# Load the medium French model (downloaded separately)
nlp = spacy.load("fr_core_news_md")

# Process a French sentence
doc = nlp("Il avait mangé la pomme hier soir.")
\end{lstlisting}

\section{Tokens: The Fundamental Unit}

When spaCy processes text, it splits it into \textbf{tokens}. A token is usually a word,
but can also be punctuation or a contracted form. The resulting \texttt{doc} object
behaves like a list of \texttt{Token} objects:

\begin{lstlisting}[caption={Inspecting spaCy tokens}]
for token in doc:
    print(f"{token.text:15} | POS: {token.pos_:8} | Lemma: {token.lemma_}")

# Output:
# Il              | POS: PRON     | Lemma: il
# avait           | POS: AUX      | Lemma: avoir
# mangé           | POS: VERB     | Lemma: manger
# la              | POS: DET      | Lemma: le
# pomme           | POS: NOUN     | Lemma: pomme
# hier            | POS: ADV      | Lemma: hier
# soir            | POS: NOUN     | Lemma: soir
# .               | POS: PUNCT    | Lemma: .
\end{lstlisting}

\section{Lemmatization: Finding the Root Form}

A \textbf{lemma} is the dictionary form of a word—the form you would look up in a
dictionary. ``mangé'', ``mange'', ``mangerons'', and ``mangions'' all share the lemma
``manger.'' Lemmatization allows us to count all forms of a verb as a single vocabulary
item.

\texttt{token.lemma\_} gives us the lemma of any token.

\section{Part-of-Speech Tagging}

\textbf{Part-of-Speech (POS) tagging} assigns each token a grammatical category:

\begin{table}[H]
\centering
\caption{spaCy POS tags used in FrenchBooks}
\begin{tabularx}{\textwidth}{l l X}
\toprule
\textbf{POS Tag} & \textbf{French Name} & \textbf{Example} \\
\midrule
VERB & Verbe & manger, courir \\
AUX  & Auxiliaire & avoir, être \\
NOUN & Nom & pomme, maison \\
ADJ  & Adjectif & grand, petit \\
ADV  & Adverbe & vite, hier \\
PRON & Pronom & il, elle, nous \\
DET  & Déterminant & le, la, un \\
PUNCT & Ponctuation & . , ! \\
\bottomrule
\end{tabularx}
\end{table}

\texttt{token.pos\_} gives us the POS tag as a string.

\section{Morphological Analysis: Extracting Verb Tenses}

\textbf{Morphological analysis} examines the internal structure of a word to extract
grammatical features. For French verbs, these features include:

\begin{itemize}
  \item \textbf{Tense:} Pres (present), Past, Imp (imperfect), Fut (future)
  \item \textbf{Mood:} Ind (indicative), Sub (subjunctive), Cnd (conditional), Imp (imperative)
  \item \textbf{Person:} 1, 2, 3
  \item \textbf{Number:} Sing, Plur
  \item \textbf{VerbForm:} Inf (infinitive), Part (participle), Fin (finite)
\end{itemize}

\begin{lstlisting}[caption={Extracting morphological features from a verb (LLR-15)}]
for token in doc:
    if token.pos_ in ("VERB", "AUX"):
        morph = token.morph.to_dict()
        tense  = morph.get("Tense",    "N/A")
        mood   = morph.get("Mood",     "N/A")
        person = morph.get("Person",   "N/A")
        number = morph.get("Number",   "N/A")
        print(
            f"{token.text} -> Tense={tense}, "
            f"Mood={mood}, Person={person}, Number={number}"
        )

# Output for "avait mangé":
# avait -> Tense=Past, Mood=Ind, Person=3, Number=Sing
# mangé -> Tense=Past, Mood=Ind, Person=N/A, Number=Sing
\end{lstlisting}

\section{French Tense Reference Chart}

\begin{table}[H]
\centering
\caption{French tenses and their spaCy morphological signatures}
\begin{tabularx}{\textwidth}{l l l}
\toprule
\textbf{French Tense} & \textbf{spaCy Tense} & \textbf{spaCy Mood} \\
\midrule
Présent de l'indicatif & Pres & Ind \\
Imparfait & Imp & Ind \\
Passé simple & Past & Ind \\
Futur simple & Fut & Ind \\
Conditionnel présent & Pres & Cnd \\
Subjonctif présent & Pres & Sub \\
Subjonctif imparfait & Imp & Sub \\
Impératif & N/A & Imp \\
\bottomrule
\end{tabularx}
\end{table}

% ============================================================
%  PART IV — THE DASHBOARD
% ============================================================
\part{The Library Dashboard}

% ============================================================
\chapter{Word Frequency Heatmaps with matplotlib}
\label{ch:heatmap}
% ============================================================

\section{What is a Word Heatmap?}

A word frequency heatmap is a visual representation of how often words appear in a text.
High-frequency words are shown in warm colors (red/orange); low-frequency words in cool
colors (blue). This allows a student to immediately see which vocabulary is central to a
book.

\section{Computing Word Frequencies}

\begin{lstlisting}[caption={Word frequency computation (LLR-17)}]
from collections import Counter

# Collect all lemmas, excluding stop words and punctuation
word_counts = Counter(
    token.lemma_
    for token in doc
    if not token.is_stop and not token.is_punct
)

# Get the 50 most common words
top_words = word_counts.most_common(50)
\end{lstlisting}

\section{Embedding matplotlib in PyQt6}

By default, calling \texttt{plt.show()} opens a separate matplotlib window. To render
charts \textit{inside} a PyQt6 widget, we use the Qt Agg backend:

\begin{lstlisting}[caption={Embedding a matplotlib figure in a Qt widget (LLR-18)}]
from matplotlib.backends.backend_qtagg import FigureCanvasQTAgg
from matplotlib.figure import Figure
import numpy as np

class HeatmapCanvas(FigureCanvasQTAgg):
    def __init__(self, word_counts, parent=None):
        fig = Figure(figsize=(10, 4))
        super().__init__(fig)
        self.setParent(parent)
        ax = fig.add_subplot(111)
        words, counts = zip(*word_counts[:30])
        data = np.array(counts).reshape(5, 6)  # Reshape for 2D heatmap
        im = ax.imshow(data, cmap="YlOrRd", aspect="auto")
        fig.colorbar(im, ax=ax)
        ax.set_title("Word Frequency Heatmap")
        fig.tight_layout()
\end{lstlisting}

% ============================================================
%  PART IV-B — DYNAMIC NLP (feature/dynamic-nlp)
% ============================================================

% ============================================================
\chapter{Corpus-Level Analysis: CorpusWorker and DashboardController}
\label{ch:corpus}
% ============================================================

\section{Per-Page vs Corpus Analysis}

The application operates at two granularities of NLP:

\begin{description}
  \item[Per-page (Reader)] The \texttt{NLPWorker} (Chapter \ref{ch:threading})
    processes only the text visible on the current page, keeping the reader
    panel always in sync with what the user sees.
  \item[Full-corpus (Dashboard)] The \texttt{CorpusWorker} processes every
    page of the book in sequence, building a global picture of vocabulary,
    verb tenses, and word frequencies.
\end{description}

\section{CorpusWorker: Processing Every Page}

\begin{lstlisting}[caption={CorpusWorker iterates all pages (HLR-01, LLR-17, LLR-19, LLR-20)}]
class CorpusWorker(QRunnable):
    def __init__(self, nlp, page_texts: list[str]):
        super().__init__()
        self.nlp        = nlp
        self.page_texts = page_texts
        self.signals    = CorpusSignals()
        self._cancel    = threading.Event()

    def run(self) -> None:
        freq = Counter()
        tense_map = defaultdict(list)
        for idx, text in enumerate(self.page_texts):
            if self._cancel.is_set():
                return
            doc = self.nlp(text)
            for token in doc:
                if not (token.is_stop or token.is_punct):
                    freq[token.lemma_.lower()] += 1
                if token.pos_ in ("VERB", "AUX"):
                    morph = token.morph.to_dict()
                    key = (morph.get("Tense"), morph.get("Mood"),
                           morph.get("Person"), morph.get("Number"))
                    tense_map[key].append(token.sent.text[:80])
            self.signals.progress.emit(idx + 1, len(self.page_texts))
        result = CorpusResult()
        result.vocab = [(l, c, "") for l, c in freq.most_common(500)]
        result.word_counts = freq
        result.tense_rows  = [(t, m, p, n, ex[0])
                               for (t,m,p,n), ex in sorted(tense_map.items())]
        self.signals.finished.emit(result)
\end{lstlisting}

\begin{notebox}
The \texttt{threading.Event} cancellation flag works because Python's
\texttt{threading} module is thread-safe for simple boolean reads. The
\texttt{QRunnable.run()} method checks \texttt{self.\_cancel.is\_set()}
at the start of each page loop — if the user closes the book or clicks
"Analyze" again, the old worker exits cleanly.
\end{notebox}

\section{The HeatmapCanvas: Matplotlib Inside Qt}

By default, calling \texttt{plt.show()} would open a separate OS window for the
matplotlib chart. We use the \textbf{Qt Agg backend} to render the figure inside
a PyQt6 widget instead (LLR-18):

\begin{lstlisting}[caption={HeatmapCanvas embedded in Qt (LLR-17, LLR-18)}]
from matplotlib.backends.backend_qtagg import FigureCanvasQTAgg
from matplotlib.figure import Figure

class HeatmapCanvas(FigureCanvasQTAgg):
    def __init__(self, word_counts: Counter, top_n: int = 40, parent=None):
        fig = Figure(figsize=(9, 3), dpi=96)
        super().__init__(fig)
        ax = fig.add_subplot(111)
        words, counts = zip(*word_counts.most_common(top_n))
        grid = np.array(counts).reshape(5, -1)
        ax.imshow(grid, cmap="YlOrRd", aspect="auto")
        fig.colorbar(ax.images[0], ax=ax, label="Frequency")
        fig.tight_layout()
        self.draw()
\end{lstlisting}

The key insight: \texttt{FigureCanvasQTAgg} is itself a \texttt{QWidget}.
You can add it to any Qt layout just like a button or a label.

\section{DashboardController Wiring}

The \texttt{DashboardController} is injected into \texttt{MainWindow.set\_controller()}
\textit{after} the spaCy model has been loaded in the background:

\begin{enumerate}
  \item \texttt{DashboardView.analyze\_requested} signal fires with the file path.
  \item \texttt{DashboardController.on\_analyze\_requested()} extracts all page texts
    (using fitz for PDFs, ebooklib+BeautifulSoup4 for EPUBs).
  \item A \texttt{CorpusWorker} is dispatched to \texttt{QThreadPool}.
  \item \texttt{CorpusSignals.finished} delivers the \texttt{CorpusResult} to
    \texttt{DashboardController.\_on\_analysis\_done()}.
  \item The controller calls \texttt{view.populate\_vocab()},
    \texttt{view.populate\_tenses()}, and \texttt{view.replace\_heatmap(canvas)}.
\end{enumerate}

% ============================================================
%  PART IV-C — CLOZE FLASHCARD ENGINE (feature/cloze-flashcard-engine)
% ============================================================

% ============================================================
\chapter{Cognitive Load and Memory Retention: Implementing Cloze Deletion}
\label{ch:cloze}
% ============================================================

\section{The Science: Why Full-Sentence Context Beats Word Lists}

Consider two ways to study the French verb \textit{s'évanouir} (to faint):

\begin{description}
  \item[Method A — Isolated pair] \hfill \\
    Front: \textit{s'évanouir} \quad Back: to faint

  \item[Method B — Cloze deletion] \hfill \\
    Front: \textit{``Elle \_\_\_\_\_ en voyant le sang.''}
    (``She \_\_\_\_\_ upon seeing the blood.'') \\
    Back: \textit{s'évanouir — to faint}
\end{description}

Method B is demonstrably superior. Several decades of cognitive-science research explain
why:

\begin{enumerate}
  \item \textbf{The Generation Effect (Slamecka \& Graf, 1978).} When learners must
    \textit{generate} the missing word from context clues rather than simply reading it,
    encoding is deeper and recall is more durable.

  \item \textbf{The Context Effect (Craik \& Tulving, 1975).} Words learned in rich
    sentence contexts are encoded with more associative cues. At retrieval time, the
    surrounding words act as additional pathways back to the target.

  \item \textbf{Input Hypothesis (Krashen, 1985).} Comprehensible input — real sentences
    from real texts that are \textit{just slightly} above the learner's current level —
    is the primary mechanism of language acquisition.

  \item \textbf{Spaced Repetition Amplification.} Anki's SM-2 spaced-repetition algorithm
    schedules reviews just before you are predicted to forget. When each card already
    encodes a sentence (not just a word), each review is simultaneously a reading
    comprehension and vocabulary exercise.
\end{enumerate}

\begin{notebox}
The term ``cloze'' comes from the psychological concept of \textbf{closure} —
the mind's tendency to complete incomplete patterns. Wilson Taylor coined the term
in 1953 for his procedure of deleting every $n$th word from a passage and asking
readers to fill in the blanks as a measure of reading comprehension.
\end{notebox}

\section{Extracting Sentences with spaCy: The \texttt{token.sent} Span}

When spaCy processes text with a pipeline that includes a dependency parser or sentence
segmenter (both included in \texttt{fr\_core\_news\_md}), it adds sentence boundary
annotations. The resulting \texttt{Doc} object is partitioned into \texttt{Span} objects
accessible via \texttt{doc.sents}.

Every token inside a sentence-annotated Doc exposes its enclosing sentence directly:

\begin{lstlisting}[caption={Extracting a sentence from a spaCy token (LLR-25)}]
import spacy

nlp = spacy.load("fr_core_news_md")
doc = nlp("Elle s'évanouit en voyant le sang. Il appela une ambulance.")

for token in doc:
    if token.text == "s'évanouit":
        sentence = token.sent        # this is a spaCy Span object
        print(sentence.text)
        # Output: "Elle s'évanouit en voyant le sang."

        cloze = sentence.text.replace(token.text, "_____", 1)
        print(cloze)
        # Output: "Elle _____ en voyant le sang."
\end{lstlisting}

\begin{warnbox}
Always check \texttt{doc.has\_annotation("SENT\_START")} before calling
\texttt{token.sent} (LLR-34). If sentence segmentation was not run (e.g.
the pipeline was loaded with \texttt{exclude=["parser"]}), accessing
\texttt{token.sent} raises a \texttt{ValueError}.
\end{warnbox}

\section{The FlashcardService}

The \texttt{FlashcardService} class in \texttt{services/flashcard\_service.py} is the
domain model for all flashcard data. Following the MVC pattern (Chapter~\ref{ch:mvc}),
it has zero PyQt6 imports — it is pure Python.

\begin{lstlisting}[caption={FlashcardRecord and FlashcardService skeleton (LLR-24, LLR-26)}]
from dataclasses import dataclass, field
from pathlib import Path
import csv

@dataclass
class FlashcardRecord:
    word:           str
    lemma:          str
    sentence:       str      # original full sentence
    cloze_sentence: str      # sentence with target word replaced by "_____"
    definition:     str = ""
    audio_path:     str = ""
    source_file:    str = ""
    page_index:     int = 0

class FlashcardService:
    def __init__(self, audio_dir: str = "flashcards/audio"):
        self._cards: list[FlashcardRecord] = []
        self._audio_dir = Path(audio_dir)
        self._audio_dir.mkdir(parents=True, exist_ok=True)

    def add_card(self, word, lemma, sentence, cloze, source_file,
                 page_index, definition="") -> FlashcardRecord:
        record = FlashcardRecord(word=word, lemma=lemma, sentence=sentence,
                                 cloze_sentence=cloze, definition=definition,
                                 source_file=source_file, page_index=page_index)
        self._cards.append(record)
        return record

    def export_csv(self, path: str) -> None:
        with open(path, "w", encoding="utf-8", newline="") as f:
            writer = csv.writer(f, delimiter=";")
            writer.writerow(["Front", "Back", "Audio"])
            for card in self._cards:
                back  = f"{card.word} — {card.definition}" if card.definition else card.word
                audio = f"[sound:{Path(card.audio_path).name}]" if card.audio_path else ""
                writer.writerow([card.cloze_sentence, back, audio])
\end{lstlisting}

\section{The Anki CSV Format Explained}

Anki's import system reads CSV files where each row is one card. We use the
\textbf{Basic (and reversed card)} note type, but the same format also works with
the \textbf{Cloze} note type. Our three columns map to Anki fields:

\begin{table}[H]
\centering
\caption{Anki CSV column mapping (LLR-27)}
\begin{tabularx}{\textwidth}{l l X}
\toprule
\textbf{CSV Column} & \textbf{Anki Field} & \textbf{Example} \\
\midrule
Front & Question & \textit{``Elle \_\_\_\_\_ en voyant le sang.''} \\
Back  & Answer   & \textit{s'évanouir — to faint} \\
Audio & Extra    & \texttt{[sound:card\_evanouit\_3f2a.mp3]} \\
\bottomrule
\end{tabularx}
\end{table}

The \texttt{[sound:filename]} syntax is Anki's media tag. When Anki encounters it in
any field, it automatically plays the referenced audio file (which must be in Anki's
\texttt{collection.media} folder).

\section{gTTS: French Audio Synthesis}

\texttt{gTTS} (Google Text-to-Speech) wraps the same TTS engine used by Google
Translate. It is synchronous, simple, and produces acceptable audio quality for French:

\begin{lstlisting}[caption={gTTS usage (LLR-29)}]
from gtts import gTTS

tts = gTTS(text="Elle s'évanouit en voyant le sang.", lang="fr", slow=False)
tts.save("output.mp3")
# output.mp3 is now a standard MP3 file playable in any audio player
\end{lstlisting}

The \texttt{slow=False} parameter uses the normal speech rate. Setting
\texttt{slow=True} produces a slower, clearer pronunciation — useful for beginners.

\begin{warnbox}
\textbf{gTTS requires an internet connection.} It sends the text to Google's servers.
Never call it on the main GUI thread — always use an \texttt{AudioWorker} (LLR-29).
If the network call fails, the \texttt{AudioWorker} catches the exception, emits an
\texttt{error} signal, and leaves the card's \texttt{audio\_path} empty (LLR-30).
\end{warnbox}

\subsection{Optional: edge-tts (Higher Quality Neural Voices)}

\texttt{edge-tts} provides access to Microsoft's Azure neural TTS voices, including
\texttt{fr-FR-DeniseNeural} — a high-quality female French voice. It is asynchronous:

\begin{lstlisting}[caption={edge-tts usage (async, LLR-35)}]
import asyncio, edge_tts

async def synthesize(text: str, output_path: str) -> None:
    comm = edge_tts.Communicate(text, voice="fr-FR-DeniseNeural")
    await comm.save(output_path)

# Inside AudioWorker.run() (background thread):
asyncio.run(synthesize(self.text, self.output_path))
\end{lstlisting}

\section{AudioController: Thread-Safe Playback}

The \texttt{AudioController} wraps Qt's \texttt{QMediaPlayer} and manages the synthesis
queue. \texttt{QMediaPlayer} is the standard Qt audio playback class:

\begin{lstlisting}[caption={AudioController playback controls (LLR-31)}]
from PyQt6.QtMultimedia import QMediaPlayer, QAudioOutput
from PyQt6.QtCore import QUrl

class AudioController(QObject):
    def __init__(self, parent=None):
        super().__init__(parent)
        self._player = QMediaPlayer()
        self._output = QAudioOutput()
        self._player.setAudioOutput(self._output)

    def play(self, file_path: str) -> None:
        self._player.setSource(QUrl.fromLocalFile(file_path))
        self._player.play()

    def set_playback_rate(self, rate: float) -> None:
        self._player.setPlaybackRate(rate)  # 0.5 = half speed, 2.0 = double
\end{lstlisting}

\section{The Flashcard Manager Panel}

The \texttt{FlashcardManagerPanel} is a \texttt{QWidget} added as a tab to the
Dashboard's \texttt{QTabWidget}. It provides the complete user-facing interface for
managing, previewing, and exporting flashcards.

\begin{figure}[H]
\centering
\fbox{\parbox{0.85\textwidth}{\centering
  \textbf{[Flashcard Manager Layout]} \\[0.5em]
  \begin{tabular}{|p{0.38\textwidth}|p{0.08\textwidth}|p{0.10\textwidth}|p{0.16\textwidth}|}
  \hline
  Cloze Sentence & Word & Lemma & Audio \\ \hline
  Elle \_\_\_\_\_ en voyant & s'évanouit & s'évanouir & Ready \\
  Il \_\_\_\_\_ la porte & ouvrit & ouvrir & Pending \\
  \hline
  \end{tabular} \\[0.5em]
  [Synthesize Audio] [Export CSV...] [Delete Selected] \\
  [Play] [Pause] [Stop] \quad Speed: [1.0x ▼]
}}
\caption{Flashcard Manager Panel layout (LLR-33)}
\end{figure}

% ============================================================
%  PART V — CERTIFICATION
% ============================================================
\part{DO-178C Certification Pipeline}

% ============================================================
\chapter{GitFlow and Traceability}
\label{ch:gitflow}
% ============================================================

\section{Why GitFlow?}

GitFlow is a branching model for Git that enforces a clean separation between
\textit{development work} (happening on feature branches), \textit{integration}
(happening on develop), and \textit{production releases} (happening on master).

For a DO-178C-compliant project, GitFlow provides natural checkpoints where
verification activities can occur before code reaches the production branch.

\section{The Branch Hierarchy}

\begin{lstlisting}[language=bash, caption={Setting up the GitFlow branch structure}]
# Start on master (the production baseline)
git checkout master

# Create the long-lived integration branch
git checkout -b develop
git push -u origin develop

# Start a new feature
git checkout -b feature/native-readers
# ... implement PDF/EPUB rendering ...
git add .
git commit -m "feat(reader): render PDF page as QPixmap [LLR-05]"
git push -u origin feature/native-readers
# ... create PR to develop via: gh pr create ...
\end{lstlisting}

\section{The Commit Format (LLR-21)}

Every commit message in this project follows this format:

\begin{center}
\texttt{<type>(<scope>): <description> [<REQ-ID>]}
\end{center}

\begin{itemize}
  \item \textbf{type:} \texttt{feat} (new feature), \texttt{fix} (bug fix),
    \texttt{docs} (documentation), \texttt{refactor}, \texttt{test}, \texttt{chore}
  \item \textbf{scope:} The component being changed (e.g., \texttt{reader},
    \texttt{nlp}, \texttt{dashboard})
  \item \textbf{REQ-ID:} The HLR or LLR that this commit satisfies
\end{itemize}

\textbf{Example:} \texttt{feat(reader): render PDF page as QPixmap [LLR-05]}

This format means that by reading the git log, a certification auditor can instantly
trace every change in the codebase back to its originating requirement.

% ============================================================
\backmatter
% ============================================================

\chapter*{Glossary}
\addcontentsline{toc}{chapter}{Glossary}

\begin{description}
  \item[Corpus] A large collection of text used for linguistic analysis.
  \item[DO-178C] Software development assurance standard for airborne software, adapted
    here for its rigorous traceability requirements.
  \item[EPUB] Electronic Publication; a ZIP-based e-book format using HTML internally.
  \item[Event Loop] The main loop of a GUI application that processes user input events.
  \item[GitFlow] A Git branching model with defined roles for master, develop, feature,
    release, and hotfix branches.
  \item[GUI] Graphical User Interface; the visual, clickable part of an application.
  \item[HLR] High-Level Requirement; a system-wide behavioral requirement.
  \item[Lemma] The base/dictionary form of a word (e.g., ``manger'' is the lemma of
    ``mangé'').
  \item[LLR] Low-Level Requirement; a detailed, implementable requirement derived from
    an HLR.
  \item[MVC] Model-View-Controller; an architectural pattern separating data, display,
    and control logic.
  \item[NLP] Natural Language Processing; using computers to understand human language.
  \item[POS] Part-of-Speech; the grammatical role of a word (noun, verb, adjective, etc.).
  \item[PyMuPDF (fitz)] A Python binding to the MuPDF PDF rendering library.
  \item[PyQt6] Python bindings for the Qt 6 GUI toolkit.
  \item[QPixmap] A Qt object representing an image optimized for display on screen.
  \item[QRunnable] A Qt class for background tasks that run in a thread pool.
  \item[QThreadPool] A Qt-managed pool of background threads.
  \item[QWebEngineView] A Qt widget embedding the Chromium web engine.
  \item[Signal/Slot] Qt's thread-safe event communication mechanism.
  \item[spaCy] An industrial-strength Python NLP library.
  \item[venv] Python virtual environment; an isolated Python installation per project.
\end{description}

\printindex

\end{document}
